% =============================================================================
% 3. fejezet — Elméleti háttér
% =============================================================================

Ez a fejezet a dolgozatban felhasznált kulcstechnológiák elméleti hátterét
építi fel, az alapoktól az alkalmazott architektúráig.
% Először a mesterséges neuron és a többrétegű perceptron fogalmát, majd a
% neurális hálók tanításának mechanizmusát ismertetjük; ezt követi a konvolúciós
% hálók, a figyelmi mechanizmus és a Vision Transformer tárgyalása. A fejezet
% második fele a dolgozat módszertanához közvetlenül kapcsolódó témákat -- a
% transzfertanulást, az osztály-egyensúlytalanság kezelését, a kiértékelési
% metrikákat, a statisztikai elemzést és a kommunikációs technológiákat --
% mutatja be. A cél nem a terület teljes körű tankönyvi feldolgozása, hanem a
% dolgozat megértéséhez szükséges fogalmak tömör, mégis önállóan érthető
% összefoglalása.

% =============================================================================
\section{A mesterséges neuron és a többrétegű perceptron}
% =============================================================================

A mesterséges neurális hálózatok koncepciója a biológiai idegrendszer működéséből
merít ihletet, annak elveit igyekszik leképezni és továbbgondolni.
Éppen ezért érdemes először áttekinteni a biológiai neurális hálózat alapjait,
mielőtt a mesterséges megfelelőjére térnénk.
Az emberi idegrendszer alapegységei az idegsejtek, amelyek szerkezetileg három
jól elkülöníthető részből állnak. A dendritek a sejttestből kiágazó nyúlványok,
amelyek elsődleges feladata a külső jelek és ingerek befogadása. 
Maga a sejttest látja el a beérkező információk integrálásának és feldolgozásának
szerepét. Az axon — amely a sejttestet a szinaptikus végződésekkel köti össze — 
kulcsfontosságú az idegsejtek közötti kommunikációban, mivel ezen keresztül
továbbítódnak az elektromos jelek. Az a mechanizmus, amelynek során egy neuron
axonja egy másik neuron dendritjével kerül kapcsolatba, és a felszabaduló
neurotranszmitterek révén megtörténik a jeltovábbítás a két sejt között, szinapszis néven ismert.

A mély tanulás alapegysége a mesterséges neuron, amely a biológiai idegsejt
erősen leegyszerűsített matematikai modellje \cite{goodfellow2016deep}. A neuron
a bemeneti jelek ($x_1, \dots, x_n$) súlyozott összegét képezi, hozzáad egy
eltolási (bias) tagot, majd az eredményt egy nemlineáris \textit{aktivációs
függvényen} ($\varphi$) vezeti át:

\begin{equation}
  y = \varphi\!\left(\sum_{i=1}^{n} w_i x_i + b\right),
  \label{eq:neuron}
\end{equation}

ahol $w_i$ a bemenetekhez rendelt súlyok, $b$ az eltolás. A súlyok és az eltolás
a tanulás során hangolódó paraméterek. Egyetlen neuron csak lineárisan
szétválasztható feladatokat képes megoldani; a kifejezőerőt a neuronok
\textit{rétegekbe} szervezése adja.

A \textit{többrétegű perceptron} (multilayer perceptron, MLP) neuronok egymást
követő rétegeiből épül fel: egy bemeneti rétegből, egy vagy több rejtett
rétegből és egy kimeneti rétegből, ahol az egyes rétegek neuronjai a következő
réteg minden neuronjához kapcsolódnak (\ref{fig:neuron-mlp}.~ábra). Egy elegendően
széles, nemlineáris aktivációval rendelkező rejtett réteggel az MLP tetszőleges
folytonos függvényt képes tetszőleges pontossággal közelíteni -- ez az
\textit{univerzális approximáció} tulajdonsága --, a gyakorlatban azonban a
mélység (a rétegek száma) teszi lehetővé a bonyolult, hierarchikus
reprezentációk hatékony tanulását \cite{goodfellow2016deep}.

\begin{figure}[H]
  \centering
  \includegraphics[width=\textwidth]{images/diagrams/neuron_mlp.png}
  \caption{Balra: a mesterséges neuron, amely a bemenetek súlyozott összegét egy
    aktivációs függvényen vezeti át. Jobbra: a többrétegű perceptron, amelyben a
    teljesen összekötött rétegek a bemenetből a kimenet felé továbbítják a jelet.}
  \label{fig:neuron-mlp}
\end{figure}

% =============================================================================
\section{Aktivációs függvények}
% =============================================================================

Az aktivációs függvény viszi be a nemlinearitást a hálóba; e nélkül a rétegek
egymásutánja egyetlen lineáris transzformációvá esne össze, és a háló
kifejezőereje nem haladná meg egyetlen réteg kifejezőerejét. A legelterjedtebb választások a
következők \cite{goodfellow2016deep}:

\begin{itemize}
  \item \textbf{Szigmoid:} $\sigma(z) = 1/(1 + e^{-z})$, amely a bemenetet a
  $(0;1)$ intervallumra képezi; bináris osztályozás kimeneti rétegében
  valószínűségként értelmezhető, rejtett rétegekben azonban a telítődő
  tartományokban eltűnő gradienst okozhat.
  \item \textbf{Tangens hiperbolikus:} $\tanh(z)$, amely a $(-1;1)$ intervallumra
  képez, és nulla középpontú kimenetével gyakran kedvezőbb a szigmoidnál.
  \item \textbf{Rektifikált lineáris egység (ReLU):} $\mathrm{ReLU}(z) = \max(0,
  z)$, amely egyszerűsége és a pozitív tartományban nem telítődő, gradiens-barát
  viselkedése miatt a mély hálók alapértelmezett választásává vált. Számos
  változata létezik, például a negatív tartományban kis meredekséget engedő
  \textit{Leaky ReLU}, illetve a sima \textit{GELU}, amelyet a Transformer-alapú
  modellek -- köztük a dolgozat ViT modellje -- is alkalmaznak.
\end{itemize}

A többosztályos kimeneti rétegben a \textbf{softmax} függvény használatos, amely
a kimeneti pontszámokat (logitokat) egy osztályok feletti valószínűségi
eloszlássá normalizálja:

\begin{equation}
  \mathrm{softmax}(z)_k = \frac{e^{z_k}}{\sum_{j} e^{z_j}}.
  \label{eq:softmax}
\end{equation}

% =============================================================================
\section{A neurális hálók tanítása}
% =============================================================================

A háló tanítása a paraméterek (súlyok és eltolások) olyan beállítását jelenti,
amely minimalizálja a modell előrejelzései és a valós címkék közötti eltérést. Ezt
az eltérést egy \textit{veszteségfüggvény} (loss function) számszerűsíti, a
minimalizálást pedig gradiens-alapú optimalizálás végzi.

\subsection{Veszteségfüggvények}

A feladat jellegétől függően különböző veszteségfüggvényeket alkalmazunk.
Regressziónál az átlagos négyzetes eltérés (mean squared error) a tipikus
választás. Osztályozásnál a \textit{keresztentrópia} (cross-entropy) használatos,
amely a becsült és a valós valószínűségi eloszlás közötti eltérést méri; bináris
esetben \cite{goodfellow2016deep}:

\begin{equation}
  \mathcal{L}_{\mathrm{BCE}} = -\frac{1}{N}\sum_{i=1}^{N}
  \big[ y_i \log p_i + (1 - y_i)\log(1 - p_i) \big],
  \label{eq:bce}
\end{equation}

ahol $y_i \in \{0,1\}$ a valós címke, $p_i$ a modell által a pozitív osztályra
becsült valószínűség. Erősen kiegyensúlyozatlan feladatban e veszteség
módosítására van szükség; ezt a \ref{sec:imbalance}.~szakasz tárgyalja.

\subsection{Gradiens-alapú tanulás és hiba-visszaterjesztés}

A veszteség minimalizálása \textit{gradiensereszkedéssel} (gradient descent)
történik: a paramétereket a veszteségfüggvény gradiensével ellentétes irányban,
kis lépésekben mozgatjuk, hiszen a gradiens a legmeredekebb növekedés irányába
mutat. Egy $\theta$ paraméter frissítése:

\begin{equation}
  \theta \leftarrow \theta - \eta \, \frac{\partial \mathcal{L}}{\partial \theta},
  \label{eq:gd}
\end{equation}

ahol $\eta$ a \textit{tanulási ráta} (learning rate), amely a lépés nagyságát
szabályozza. A többrétegű hálók esetén a gradienseket a \textit{hiba-visszaterjesztés}
(backpropagation) algoritmusa számítja ki, amely a deriválás láncszabályát
alkalmazva a kimeneti hibát rétegről rétegre, visszafelé terjeszti, és így
minden paraméterre hatékonyan meghatározza a gradienst \cite{goodfellow2016deep}.

\subsection{Optimalizálók: az SGD-től az AdamW-ig}

A teljes adathalmaz feletti gradiens kiszámítása nagy adathalmazokon költséges,
ezért a gyakorlatban a \textit{sztochasztikus gradiensereszkedést} (stochastic
gradient descent, SGD) alkalmazzák: a gradienst a tanítóadat kisebb,
véletlenszerű részhalmazain (\textit{batch}-ein) becsülik. Az alapváltozatot
számos technika javítja, például a \textit{momentum}, amely a korábbi lépések
irányát is figyelembe veszi, gyorsítva a konvergenciát.

A modern mély tanulás egyik legelterjedtebb optimalizálója az \textbf{Adam},
amely a gradiensek első és második momentumának mozgóátlagát tartja karban, és
ezek alapján paraméterenként adaptálja a tanulási rátát \cite{kingma2015adam}. Az
Adam egy továbbfejlesztett változata az \textbf{AdamW}, amely a súlycsökkentést
(weight decay) -- egy, a túltanulást mérséklő regularizációs tagot -- a gradiens
frissítésétől szétválasztva, helyesebben kezeli, és így jobb általánosítást ér el
\cite{loshchilov2019adamw}. A dolgozat valamennyi modelljét AdamW
optimalizálóval tanítottuk.

A tanulási ráta időbeli alakítása (\textit{ütemezés}, scheduling) szintén
befolyásolja a konvergenciát. Bevett gyakorlat a kezdeti néhány epochban
alkalmazott \textit{bemelegítés} (warmup), amely fokozatosan emeli a tanulási
rátát, megelőzve a kezdeti nagy, destabilizáló lépéseket, majd egy
\textit{koszinusz} ütemezés, amely a ráta fokozatos csökkentésével finomhangolja
a modellt.

\subsection{Regularizáció és a túltanulás megelőzése}

A nagy kapacitású hálók hajlamosak a \textit{túltanulásra} (overfitting): a
tanítóadatot kívülről megtanulják, de új adatokra rosszul általánosítanak. Ennek
megelőzésére több technika áll rendelkezésre:

\begin{itemize}
  \item \textbf{Dropout:} a tanítás során a neuronok egy véletlenszerűen
  kiválasztott hányadát ideiglenesen kikapcsolja, így a háló nem támaszkodhat
  egyetlen jellemzőre sem, ami robusztusabb reprezentációkat eredményez
  \cite{srivastava2014dropout}.
  \item \textbf{Súlycsökkentés (weight decay):} a nagy súlyokat bünteti, a
  modellt egyszerűbb megoldások felé terelve.
  \item \textbf{Adatkiterjesztés (augmentation):} a tanítóképek véletlenszerű
  transzformációival (tükrözés, forgatás, szín-jitter) mesterségesen növeli az
  adathalmaz változatosságát; ezt a dolgozat módszertani része részletezi.
  \item \textbf{Korai leállítás (early stopping):} a tanítást akkor állítja le,
  amikor a validációs teljesítmény már nem javul, megelőzve a felesleges
  túltanulást.
\end{itemize}

A tanítás stabilitását és gyorsaságát a \textbf{batch-normalizálás} is segíti,
amely a rétegek aktivációit batch-enként normalizálja, csökkentve a belső
eloszlás-eltolódást és lehetővé téve a nagyobb tanulási ráták használatát
\cite{ioffe2015batchnorm}.

% =============================================================================
\section{Konvolúciós neurális hálók}
% =============================================================================

A teljesen összekötött MLP nem skálázódik jól képekre: egy közepes méretű kép
több százezer képpontból áll, így a teljes összeköttetés kezelhetetlenül sok
paramétert eredményezne, és figyelmen kívül hagyná a képek térbeli szerkezetét. A
\textit{konvolúciós neurális hálók} (CNN-ek) erre adnak választ a súlymegosztás
és a lokális kapcsolatok elvével \cite{goodfellow2016deep}.

\subsection{A konvolúciós réteg}

A konvolúciós réteg kis méretű, tanulható szűrőket (\textit{kernel}eket) csúsztat
végig a képen, és minden pozícióban a szűrő alatti képrészlet és a szűrő közötti
súlyozott összeget számítja ki. Egy-egy szűrő a teljes képen ugyanazokat a
súlyokat használja (súlymegosztás), így egy adott lokális mintázatot (élt,
sarkot, textúrát) a kép bármely pontján fel tud ismerni, miközben a paraméterek
száma alacsony marad. Egy réteg több szűrője különböző jellemzőkre érzékeny
\textit{jellemzőtérképeket} (feature maps) állít elő. A rétegek egymásra
építésével a háló egyre összetettebb, hierarchikus jellemzőket tanul: az alsóbb
rétegek élekre és textúrákra, a felsőbbek alakzatokra és objektumrészekre
reagálnak.

\subsection{Az összevonó réteg}

A konvolúciós rétegeket gyakran \textit{összevonó} (pooling) rétegek követik,
amelyek a jellemzőtérképek térbeli felbontását csökkentik: egy kis ablakon belül
a maximumot (max pooling) vagy az átlagot (average pooling) veszik. Ez egyrészt
csökkenti a számítási költséget és a paraméterszámot, másrészt bizonyos fokú
eltolás-invarianciát biztosít. A jellemzőkinyerés (konvolúció + összevonás
ismételt alkalmazása) után egy vagy több teljesen összekötött réteg végzi az
osztályozást (\ref{fig:cnn}.~ábra).

\begin{figure}[H]
  \centering
  \includegraphics[width=\textwidth]{images/diagrams/cnn_schematic.png}
  \caption{A konvolúciós neurális háló sematikus felépítése: a konvolúciós és
    összevonó rétegek ismételt alkalmazásával egyre absztraktabb
    jellemzőtérképeket nyer ki a képből, majd egy teljesen összekötött réteg végzi
    az osztályozást.}
  \label{fig:cnn}
\end{figure}

\subsection{Reziduális kapcsolatok és a ResNet}

A hálók mélyítése elvileg növeli a kifejezőerőt, a nagyon mély hálók tanítása
azonban az eltűnő gradiens jelensége miatt instabillá válhat. A ResNet
architektúra \cite{he2016resnet} ezt \textit{reziduális kapcsolatokkal}
(residual/skip connections) oldja meg: az egyes rétegblokkok nem közvetlenül a
kívánt leképezést, hanem ahhoz képest egy maradékot (reziduumot) tanulnak, és a
blokk bemenetét egy rövidre záró kapcsolat hozzáadja a kimenethez. Ez stabilizálja
a gradiensek áramlását, és lehetővé teszi a több tucat vagy száz réteg mélységű
hálók megbízható tanítását. A dolgozatban a ResNet-18 -- egy 18 réteg mélységű,
mintegy 11~millió paraméteres reziduális háló -- szolgál konvolúciós viszonyítási
alapként.

% =============================================================================
\section{A figyelmi mechanizmus és a Transformer}
% =============================================================================

A Transformer-architektúra alapja az úgynevezett \textit{skálázott pontszorzatos
figyelem} (scaled dot-product attention), amely a bemenet elemei közötti
összefüggéseket azok távolságától függetlenül modellezi \cite{vaswani2017attention}.
A mechanizmus minden bemeneti elemhez három vektort rendel: egy lekérdezést
($Q$, query), egy kulcsot ($K$, key) és egy értéket ($V$, value). A kimenet az
értékvektorok súlyozott összege, ahol a súlyokat a lekérdezések és kulcsok
hasonlósága határozza meg:

\begin{equation}
  \mathrm{Attention}(Q, K, V) = \mathrm{softmax}\!\left(\frac{QK^{\top}}{\sqrt{d_k}}\right) V,
  \label{eq:attention}
\end{equation}

ahol $d_k$ a kulcsvektorok dimenziója; a $\sqrt{d_k}$ skálázó tényező a
pontszorzatok nagyságát stabilizálja, megakadályozva a softmax telítődését. A
gyakorlatban a modell több, párhuzamos figyelmi „fejet'' (multi-head attention)
alkalmaz, amelyek különböző reprezentációs alterekben tanulnak összefüggéseket,
így gazdagabb leírást képesek megragadni \cite{vaswani2017attention}.

A Transformer-kódoló e figyelmi rétegeket helyzetenkénti, teljesen összekötött
előrecsatolt rétegekkel (feed-forward), maradék (reziduális) kapcsolatokkal és
rétegnormalizálással kombinálja. Mivel a figyelem önmagában érzéketlen az elemek
sorrendjére, a bemenethez \textit{helyzetkódolást} (positional encoding) adnak,
amely a sorrendi információt visszacsempészi a modellbe.

% =============================================================================
\section{A Vision Transformer architektúrája}
% =============================================================================

A Vision Transformer (ViT) a természetesnyelv-feldolgozásban bevált
Transformer-kódolót ülteti át a képosztályozásra \cite{dosovitskiy2021vit}. A
feldolgozás lépései a következők:

\begin{enumerate}
  \item \textbf{Foltokra bontás és beágyazás.} A $H \times W$ méretű bemeneti
  képet rögzített méretű, nem átfedő foltokra (jellemzően $16 \times 16$ pixel)
  bontja. Minden foltot vektorrá lapít, majd egy lineáris vetítés egy rögzített
  dimenziójú beágyazási térbe képezi -- ezek a „képtokenek'', a nyelvi modellek
  szótokenjeinek analógiái.
  \item \textbf{Osztályozó token és helyzetkódolás.} A foltbeágyazások elé egy
  tanulható osztályozó tokent (\texttt{[CLS]}) fűz, amely a teljes kép összegzett
  reprezentációját hivatott hordozni. A sorrendi (térbeli) információt tanulható
  helyzetkódolás biztosítja.
  \item \textbf{Transformer-kódoló.} A tokensorozatot több, egymásra épülő
  Transformer-kódoló réteg dolgozza fel; a többfejű önfigyelem révén minden folt
  közvetlen kapcsolatba kerülhet az összes többivel, így a modell a kép globális
  összefüggéseit is megragadja.
  \item \textbf{Osztályozó fej.} A végső réteg \texttt{[CLS]} tokenjének
  reprezentációjára egy osztályozó fej (jellemzően egy- vagy kétrétegű
  perceptron) épül, amely az osztályok feletti döntést szolgáltatja.
\end{enumerate}

A konvolúciós hálókkal szemben a ViT nem rendelkezik beépített képi induktív
torzítással (lokalitás, eltolás-invariancia), ezért nagyobb adatmennyiséget vagy
erős előtanítást igényel; cserébe a globális összefüggéseket rugalmasabban
modellezi, és nagy adathalmazokon a CNN-ekkel versenyképes vagy azokat meghaladó
teljesítményt ér el \cite{dosovitskiy2021vit}. A dolgozat fő modellje egy
ImageNet-en előtanított \texttt{vit\_base\_patch16\_224} architektúra, amelynek
sematikus felépítését a \ref{fig:vit-arch}.~ábra mutatja.

\begin{figure}[H]
  \centering
  \includegraphics[width=\textwidth]{images/diagrams/vit_architecture.png}
  \caption{A Vision Transformer (ViT) sematikus felépítése: a kép foltokra
    bontása, lineáris beágyazás és helyzetkódolás, a \texttt{[CLS]} tokennel
    kiegészített tokenszekvencia feldolgozása a Transformer-kódoló rétegeiben,
    végül a \texttt{[CLS]} reprezentációra épülő osztályozó fej, amely a
    jó\-indulatú/rosszindulatú döntést adja.}
  \label{fig:vit-arch}
\end{figure}

% =============================================================================
\section{Transzfertanulás és finomhangolás}
% =============================================================================

Az orvosi képadatok jellemzően korlátozott mennyiségben állnak rendelkezésre,
ami megnehezíti a nagy kapacitású modellek -- különösen a ViT -- nulláról történő
betanítását. A \textit{transzfertanulás} (transfer learning) erre kínál
megoldást: a modellt először egy nagy, általános képadatbázison (jellemzően az
ImageNet adathalmazon \cite{russakovsky2015imagenet}) tanítják be, ahol általános
vizuális reprezentációkat -- éleket, textúrákat, alakzatokat -- sajátít el. Az
így előtanított modellt ezután a célfeladat kisebb adathalmazán
\textit{finomhangolják} (fine-tuning). Raghu és munkatársai részletesen
elemezték a transzfertanulás szerepét az orvosi képfeldolgozásban, és rámutattak,
hogy az előtanított reprezentációk elsősorban a konvergencia gyorsításában és a
kis adathalmazokon való stabil tanulásban segítenek \cite{raghu2019transfusion}.

A finomhangolás során bevett gyakorlat a \textit{differenciált tanulási ráta}
alkalmazása: az előtanított törzshálót (backbone) kis tanulási rátával frissítik,
hogy a megtanult reprezentációk ne sérüljenek, míg az újonnan inicializált
osztályozó fejet nagyobb rátával tanítják, mivel azt a feladatra nulláról kell
megtanulni. A tanulás stabilitását -- a már említett -- bemelegítés is segíti: a kezdeti
néhány epochban fokozatosan emeli a tanulási rátát. Erre azért van szükség, mert
a tanítás elején a még betanítatlan osztályozó fej nagy hibát, és ezáltal nagy
gradienseket termel; magas tanulási ráta mellett ezek nagy súlymódosító
lépéseket okoznának, amelyek elronthatnák az előtanított törzsháló gondosan
kialakult reprezentációit.

% =============================================================================
\section{Az osztály-egyensúlytalanság kezelése}
\label{sec:imbalance}
% =============================================================================

Az orvosi szűrési feladatok jellemzően erősen kiegyensúlyozatlanok: a pozitív
(beteg) esetek a teljes adathalmaznak csupán töredékét teszik ki. Ilyenkor a
hagyományos keresztentrópia-veszteséget a sok, könnyen osztályozható negatív
példa uralja, és a modell hajlamos a többségi osztály felé torzítani. A
szakirodalom több, egymást kiegészítő technikát kínál e probléma kezelésére,
amelyeket a dolgozat kombináltan alkalmaz.

\paragraph{Súlyozott mintavételezés.} Az egyik megközelítés a tanító batch-ek
osztályeloszlásának kiegyensúlyozása: minden mintát az osztálya gyakoriságával
fordítottan arányos súllyal választunk, így a ritka pozitív esetek a tanítás
során felülreprezentálttá válnak.

\paragraph{Fókuszveszteség (focal loss).} A Lin és munkatársai által bevezetett
fókuszveszteség a keresztentrópiát úgy alakítja át, hogy a jól osztályozott
példák hozzájárulását lecsökkenti, és a tanulást a nehéz, félreosztályozott
esetekre összpontosítja \cite{lin2017focal}:

\begin{equation}
  \mathrm{FL}(p_t) = -\,\alpha_t \,(1 - p_t)^{\gamma}\, \log(p_t),
  \label{eq:focal}
\end{equation}

ahol $p_t$ a helyes osztályhoz rendelt becsült valószínűség, $\gamma \ge 0$ a
fókuszáló paraméter (amely a könnyű példák súlyát csökkenti), $\alpha_t$ pedig
egy osztálysúly, amely a pozitív és negatív esetek közötti egyensúlyt
finomhangolja. A $\gamma = 0$ választás visszaadja a súlyozott
keresztentrópiát.

\paragraph{Kétfázisú tanítás.} A kiegyensúlyozott mintavételezés mellett betanított
modell kimeneti valószínűségei eltérhetnek a valós, erősen kiegyensúlyozatlan
előfordulási aránytól. Ezt egy második, kalibráló fázis korrigálhatja, amelyben
a modellt a valós eloszlású adaton, kisebb tanulási rátával tanítjuk tovább, így
a kimeneti valószínűségek a tényleges priorhoz igazodnak (\ref{fig:two-phase}.~ábra).

\begin{figure}[H]
  \centering
  \includegraphics[width=\textwidth]{images/diagrams/two_phase_training.png}
  \caption{A kétfázisú tanítási stratégia: az első, kiegyensúlyozott fázis a
    diszkriminatív jellemzők megtanulását célozza, a második, kalibráló fázis
    pedig a kimeneti valószínűségeket a valós előfordulási arányhoz igazítja.}
  \label{fig:two-phase}
\end{figure}

% =============================================================================
\section{Kiértékelési metrikák}
% =============================================================================

Erősen kiegyensúlyozatlan feladatban a nyers pontosság (accuracy) félrevezető,
hiszen a többségi osztály triviális megjósolásával is magas érték érhető el.
Ezért orvosi szűrésnél a tévesztési mátrixból (confusion matrix) származtatott,
osztályérzékeny metrikákat alkalmazunk. Jelölje $TP$, $TN$, $FP$ és $FN$ rendre a
valódi pozitív, valódi negatív, hamis pozitív és hamis negatív esetek számát.

\begin{align}
  \text{Érzékenység (sensitivity)} &= \frac{TP}{TP + FN}, &
  \text{Specificitás (specificity)} &= \frac{TN}{TN + FP}, \label{eq:sens-spec}\\
  \text{PPV} &= \frac{TP}{TP + FP}, &
  \text{NPV} &= \frac{TN}{TN + FN}. \label{eq:ppv-npv}
\end{align}

Az \textbf{érzékenység} azt méri, hogy a modell a valóban rosszindulatú eseteknek
mekkora hányadát ismeri fel -- szűrőeszköznél ez a legkritikusabb mutató, mivel a
fel nem ismert rák (hamis negatív) a legsúlyosabb kimenetel. A \textbf{negatív
prediktív érték} (NPV) azt fejezi ki, hogy a „jóindulatú'' besorolás mennyire
megbízható.

\paragraph{ROC-görbe és AUC.} A bináris osztályozók döntési küszöbtől független
értékelésére a \textit{működési jelleggörbét} (Receiver Operating Characteristic,
ROC) használjuk, amely a valódi pozitív arányt (érzékenység) a hamis pozitív
arány ($FPR = FP/(FP+TN)$) függvényében ábrázolja a küszöb változtatásával. A
görbe alatti terület (Area Under the Curve, AUC) egyetlen számba sűríti a
rangsorolási képességet: értéke 0,5 (véletlen találgatás) és 1,0 (tökéletes
elkülönítés) között mozog.

\paragraph{Részleges AUC (pAUC).} Erős egyensúlytalanságnál és klinikai
kontextusban gyakran csak a görbe egy meghatározott szakasza érdekes. Az ISIC
2024 verseny elsődleges metrikája a 80\%-os érzékenység feletti tartományra
korlátozott részleges AUC (partial AUC), amely kifejezetten a klinikailag
elvárt, magas érzékenységi tartományban méri a teljesítményt; a metrika
értéke $[0;\ 0{,}2]$ tartományba esik, ahol $0{,}2$ a tökéletes osztályozó
értéke \cite{isic2024}.

\paragraph{Kalibráció és Brier-pontszám.} A jó rangsorolás (magas AUC) nem
garantálja, hogy a modell által kiadott valószínűségek megbízhatóak. A
\textit{kalibráció} azt méri, hogy a becsült valószínűségek mennyire felelnek
meg a tényleges bekövetkezési aránynak; Guo és munkatársai rámutattak, hogy a
modern neurális hálók gyakran rosszul kalibráltak \cite{guo2017calibration}. A
valószínűségi kimenetek minőségének egyik összegző mérőszáma a Brier-pontszám, a becsült valószínűség és a
tényleges címke közötti átlagos négyzetes eltérés:

\begin{equation}
  \mathrm{Brier} = \frac{1}{N} \sum_{i=1}^{N} (p_i - y_i)^2,
  \label{eq:brier}
\end{equation}

ahol $p_i$ a $i$-edik mintára becsült valószínűség, $y_i \in \{0,1\}$ a valódi
címke; az alacsonyabb érték jobb kalibrációt jelez.

% =============================================================================
\section{Statisztikai elemzés}
% =============================================================================

Mivel a teszthalmazban a pozitív esetek száma korlátozott, a pontbecslések
(érzékenység, specificitás, AUC) mellett azok bizonytalanságát is jelölni kell.
Az arányjellegű metrikákhoz (érzékenység, specificitás) az egzakt, binomiális
eloszláson alapuló Clopper--Pearson-féle konfidencia-intervallum ad megbízható
becslést kis mintaszám esetén is. Az AUC esetében a DeLong-féle, nemparaméteres
módszer szolgáltat konfidencia-intervallumot, illetve teszi lehetővé két modell
AUC-értékének páros összehasonlítását; a módszer egy gyors, korszerű
implementációját Sun és Xu adta meg \cite{sun2014delong}. E statisztikai eszközök
biztosítják, hogy a modellek közötti különbségekről tett állítások ne csupán a
mintavételi zaj következményei legyenek.

% =============================================================================
\section{Kommunikációs technológiák: SSE és az MCP}
% =============================================================================

A dolgozatban bemutatott rendszer valós idejű, token-szintű válaszközvetítést és
egy nyelvi modellt kiszolgáló komponensekkel való kommunikációt igényel; ezt két
technológia teszi lehetővé.

\paragraph{Szerver-küldött események (SSE).} A szerver-küldött események
(Server-Sent Events) a HTML szabvány részét képező, egyirányú, szerverről
kliensre irányuló adatfolyam-mechanizmus egyetlen, tartós HTTP-kapcsolaton
keresztül \cite{whatwg_sse}. Az SSE különösen alkalmas a nagy nyelvi modellek
token-szintű, fokozatos válaszainak továbbítására, mivel a kliens a szöveget
generálással egyidejűleg, darabonként jeleníti meg, jelentősen javítva a
felhasználói élményt.

\paragraph{Model Context Protocol (MCP).} A Model Context Protocol az Anthropic
által 2024-ben közzétett nyílt szabvány, amely egységesíti a nagy nyelvi
modellek és a külső eszközök, adatforrások közötti kommunikációt \cite{mcp2024}.
Az MCP a JSON-RPC 2.0 üzenetformátumra épül \cite{jsonrpc2013}, és többféle
átviteli réteget támogat, többek között a helyi folyamatok közötti szabványos
be- és kimeneti (stdio) csatornát. A dolgozat rendszerében az MCP biztosítja a
backend és a modellt kiszolgáló komponens közötti, jól strukturált
kommunikációt; e megoldás részletes tárgyalása a mérnöki részben található.
